\chapter{Appendix}

This appendix presents additional academic, research, innovation,
social-impact, and project-related contributions made by the students
during the development of the major project.

% =========================================================
% A. PUBLICATION
% =========================================================

\section{Publication}

A research paper based on this project has been prepared for
submission.

\subsection{Publication Details}

\begin{itemize}
    \item \textbf{Title of the paper:} Mechanical Verification of
    Reasoning--Code Alignment in LLM--Driven AutoML Pipelines
    \item \textbf{Authors:} Kinnshuk Bhaduri, Samarth Adsare,
    Ojas Godambe, and Sarvesh Rakte
    \item \textbf{Journal / Conference:} (to be filled in on
    acceptance)
    \item \textbf{Publisher:} (to be filled in on acceptance)
    \item \textbf{Date of publication / acceptance:} (to be filled
    in on acceptance)
    \item \textbf{DOI / Paper ID / URL:} (to be filled in on
    acceptance)
\end{itemize}

\subsection{Publication Evidence}

The following items will be attached as evidence:
\begin{itemize}
    \item Final submitted / published paper (PDF).
    \item Acceptance letter, if applicable.
    \item Conference certificate, if applicable.
    \item DOI page, if applicable.
    \item Conference / journal proceedings page, if applicable.
\end{itemize}


% =========================================================
% B. PATENT / COPYRIGHT
% =========================================================

\section{Patent / Copyright}

Provide details of any patent application, granted patent,
copyright registration, or other intellectual property generated
from the project.

\subsection{Intellectual Property Details}

\begin{itemize}
    \item Title:
    \item Type: Patent / Copyright
    \item Applicant(s):
    \item Inventor(s) / Author(s):
    \item Application / Registration Number:
    \item Date of Filing / Registration:
    \item Status:
\end{itemize}

\subsection{Evidence}

Attach relevant patent or copyright documentation.


% =========================================================
% C. OPEN SOURCE CONTRIBUTION
% =========================================================

\section{Open Source Contribution}

The project itself is released as open--source software. The complete
stack is split across nine repositories under the project's GitHub
organisation, each dual--licensed under the permissive
\textbf{MIT OR Apache--2.0} terms (declared via an SPDX
\texttt{SPDX-License-Identifier: MIT OR Apache-2.0} header in the
source files).

\begin{itemize}
    \item \textbf{Project:} LLM--Driven AutoML with reasoning--code
    alignment verification.
    \item \textbf{Repositories:} \texttt{automl-reusables} (shared
    library), \texttt{automl-data-service},
    \texttt{automl-metafeatures-service},
    \texttt{automl-generation-service},
    \texttt{automl-analysis-service}, \texttt{automl-gateway},
    \texttt{automl-frontend} and \texttt{automl-openui} (front--ends),
    and \texttt{automl-infra} (orchestration).
    \item \textbf{Platform:} GitHub.
    \item \textbf{Student Contributors:} Samarth Adsare, Kinnshuk
    Bhaduri, Ojas Godambe, Sarvesh Rakte.
    \item \textbf{Licence:} MIT OR Apache--2.0 (dual).
\end{itemize}

\subsection{Evidence}

Repository URLs, the SPDX licence headers, and the commit history
across the seven repositories serve as evidence. (A DOI for an
archived release is pending.)


% =========================================================
% D. DATA CONTRIBUTION
% =========================================================

\section{Data Contribution}

Describe any dataset created, curated, annotated, cleaned,
published, or contributed by the students.

\begin{itemize}
    \item Dataset Name:
    \item Dataset Description:
    \item Source:
    \item Number of Records / Samples:
    \item Data Collection Method:
    \item Annotation / Preprocessing Performed:
    \item Dataset Repository / URL:
\end{itemize}

\subsection{Evidence}

Include dataset documentation, repository screenshots,
data cards, or relevant links.


% =========================================================
% E. RESEARCH GRANT PROPOSAL
% =========================================================

\section{Research Grant Proposal}

Mention any research grant proposal prepared or submitted
based on the project.

\begin{itemize}
    \item Grant / Funding Scheme:
    \item Proposal Title:
    \item Student Contributor(s):
    \item Faculty Mentor:
    \item Funding Agency:
    \item Proposed Amount:
    \item Submission Date:
    \item Status:
\end{itemize}

\subsection{Evidence}

Attach proposal documents, acknowledgement of submission,
or other relevant evidence.


% =========================================================
% F. RESEARCH COMMERCIALIZATION EFFORTS
% =========================================================

\section{Research Commercialization Efforts}

Describe efforts undertaken to translate the project or research
outcomes into practical or commercial applications.

\subsection{Startup Formation}

If applicable, describe the steps taken towards startup formation.

\begin{itemize}
    \item Startup / Venture Name:
    \item Student Founders:
    \item Problem Addressed:
    \item Proposed Solution:
    \item Current Stage:
\end{itemize}

\subsection{Technology Transfer}

Describe any technology transfer activity associated with the project.

\begin{itemize}
    \item Technology / Solution:
    \item Receiving Organization:
    \item Nature of Transfer:
    \item Date:
    \item Current Status:
\end{itemize}

\subsection{Industry Collaboration}

Describe any collaboration with an industry partner.

\begin{itemize}
    \item Industry Partner:
    \item Nature of Collaboration:
    \item Student Contributors:
    \item Project / Activity:
    \item Duration:
    \item Outcome:
\end{itemize}

\subsection{Evidence}

Attach MoUs, collaboration letters, technology-transfer documents,
industry letters, meeting records, or other relevant evidence.


% =========================================================
% G. RESEARCH CONCLAVE
% =========================================================

\section{Research Conclave}

Mention participation, presentation, demonstration, poster presentation,
or other contribution made by the students at a Research Conclave.

\begin{itemize}
    \item Event Name:
    \item Organizing Institution:
    \item Date:
    \item Project / Research Title:
    \item Student Participant(s):
    \item Type of Participation:
    \item Award / Recognition, if any:
\end{itemize}

\subsection{Evidence}

Attach participation certificates, photographs, presentation certificates,
awards, or other supporting documents.


% =========================================================
% H. SDG GOALS ACHIEVED
% =========================================================

\section{SDG Goals Achieved}

Describe the United Nations Sustainable Development Goals (SDGs)
addressed by the project.

\begin{table}[H]
\centering
\caption{SDG Goals Addressed by the Project}
\label{tab:sdg}
\begin{tabular}{p{3cm} p{4cm} p{7cm}}
\toprule
\textbf{SDG} & \textbf{Goal} & \textbf{Contribution of the Project} \\
\midrule
SDG 1 & & \\
SDG 2 & & \\
SDG 3 & & \\
SDG 4 & & \\
SDG 5 & & \\
SDG 6 & & \\
SDG 7 & & \\
SDG 8 & & \\
SDG 9 & & \\
SDG 10 & & \\
SDG 11 & & \\
SDG 12 & & \\
SDG 13 & & \\
SDG 14 & & \\
SDG 15 & & \\
SDG 16 & & \\
SDG 17 & & \\
\bottomrule
\end{tabular}
\end{table}

Students should retain only the SDGs that are genuinely relevant
to the project and provide a clear justification for each selected goal.


% =========================================================
% 9. COURSE OUTCOMES, PO AND PSO MAPPING
% =========================================================

\section{Course Outcomes, PO and PSO Mapping}

This section presents the Course Outcomes (COs), Program Outcomes
(POs), Programme Specific Outcomes (PSOs), and their mapping with
the major project.

% =========================================================
% 9.1 COURSE OUTCOMES
% =========================================================

\subsection{Course Outcomes}

The Course Outcomes (COs) of the Major Project are as follows:

\begin{enumerate}[label=\textbf{CO\arabic*:}, leftmargin=2.2cm]

    \item Articulate the nature of science, differentiate between
    objective and subjective thinking, and situate scientific inquiry
    within its historical and philosophical context.

    \item Apply inductive, deductive, and syllogistic reasoning to
    construct and evaluate valid scientific arguments and identify
    causal relationships.

    \item Apply statistical concepts including the Central Limit
    Theorem, error propagation, confidence intervals, and data
    visualization techniques to interpret experimental results.

    \item Design and execute hypothesis testing procedures using
    Z-test, T-test, and Chi-Square tests, and critically evaluate
    the validity of scientific claims.

    \item Construct mathematical models, perform dimensional analysis,
    and derive functional relationships to support theoretical research.

    \item Produce well-structured scientific documents, deliver
    conference presentations, and demonstrate ethical conduct in
    research and publication.

\end{enumerate}


% =========================================================
% 9.2 PO ATTAINMENT / PROJECT CONTRIBUTION
% =========================================================

\subsection{Project Contribution to Program Outcomes}

The major project provides opportunities for students to demonstrate
the following Program Outcomes through problem identification,
literature review, system design, implementation, experimentation,
analysis, documentation, presentation, teamwork, and project
management.

\begin{table}[H]
\centering
\caption{Project Contribution to Program Outcomes}
\label{tab:project_po_contribution}
\small
\renewcommand{\arraystretch}{1.5}

\begin{tabular}{p{1.5cm} p{4cm} p{8cm}}
\toprule
\textbf{PO} & \textbf{Outcome} & \textbf{Evidence from Project} \\
\midrule

PO1 & Engineering Knowledge &
Describe how mathematical, scientific, computing, and engineering
knowledge is applied in the project. \\[0.5cm]

PO2 & Problem Analysis &
Describe problem identification, literature review, analysis, and
formulation performed during the project. \\[0.5cm]

PO3 & Design/Development of Solutions &
Describe the design and development of the proposed solution. \\[0.5cm]

PO4 & Conduct Investigations &
Describe experimentation, modelling, testing, data analysis, and
interpretation performed. \\[0.5cm]

PO5 & Engineering Tool Usage &
Describe the software, frameworks, computational tools, and
technologies used. \\[0.5cm]

PO6 & The Engineer and The World &
Describe societal, environmental, sustainability, health, safety,
legal, or cultural considerations. \\[0.5cm]

PO7 & Ethics &
Describe ethical considerations, responsible data usage, privacy,
intellectual property, and compliance. \\[0.5cm]

PO8 & Individual and Collaborative Team Work &
Describe individual responsibilities, teamwork, coordination,
and collaboration. \\[0.5cm]

PO9 & Communication &
Describe technical documentation, presentations, demonstrations,
and communication of project outcomes. \\[0.5cm]

PO10 & Project Management and Finance &
Describe project planning, scheduling, resource management,
budgeting, and decision-making. \\[0.5cm]

PO11 & Life-Long Learning &
Describe new technologies, independent learning, research,
adaptation, and continuous skill development. \\[0.5cm]

\bottomrule
\end{tabular}
\end{table}


% =========================================================
% 9.3 PSO CONTRIBUTION
% =========================================================

\subsection{Project Contribution to Programme Specific Outcomes}

\begin{table}[H]
\centering
\caption{Project Contribution to Programme Specific Outcomes}
\label{tab:project_pso_contribution}
\renewcommand{\arraystretch}{1.5}

\begin{tabular}{p{2cm} p{4cm} p{8cm}}
\toprule
\textbf{PSO} & \textbf{Outcome} & \textbf{Evidence from Project} \\
\midrule

PSO1 &
Professional Skills &
Describe the application of mathematics, computational methods,
statistics, artificial intelligence, and data science concepts. \\[0.7cm]

PSO2 &
Successful Career &
Describe the design and implementation of the proposed solution
using relevant state-of-the-art AI and Data Science techniques. \\[0.7cm]

\bottomrule
\end{tabular}
\end{table}


% =========================================================
% CO-PO MAPPING
% =========================================================

\begin{landscape}

\section*{CO--PO Mapping}

\begin{center}
\textbf{Mapping Scale:}
3 -- High Contribution,\quad
2 -- Moderate Contribution,\quad
1 -- Low Contribution,\quad
-- -- No Significant Contribution
\end{center}

\vspace{0.5cm}

\begin{table}[H]
\centering
\caption{CO--PO Mapping}

\renewcommand{\arraystretch}{1.6}
\setlength{\tabcolsep}{9pt}

\begin{tabular}{|c|*{11}{c|}}
\hline
\textbf{CO / PO}
& \textbf{PO1}
& \textbf{PO2}
& \textbf{PO3}
& \textbf{PO4}
& \textbf{PO5}
& \textbf{PO6}
& \textbf{PO7}
& \textbf{PO8}
& \textbf{PO9}
& \textbf{PO10}
& \textbf{PO11}
\\
\hline

\textbf{CO1} & & & & & & & & & & & \\ \hline
\textbf{CO2} & & & & & & & & & & & \\ \hline
\textbf{CO3} & & & & & & & & & & & \\ \hline
\textbf{CO4} & & & & & & & & & & & \\ \hline
\textbf{CO5} & & & & & & & & & & & \\ \hline
\textbf{CO6} & & & & & & & & & & & \\ \hline

\end{tabular}
\end{table}

\end{landscape}
% =========================================================
% CO-PSO MAPPING
% =========================================================

\section*{CO--PSO Mapping}

\begin{center}
\textbf{Mapping Scale:}
3 -- High Contribution,\quad
2 -- Moderate Contribution,\quad
1 -- Low Contribution,\quad
-- -- No Significant Contribution
\end{center}

\vspace{0.3cm}

\begin{table}[H]
\centering
\caption{CO--PSO Mapping}
\renewcommand{\arraystretch}{1.5}
\setlength{\tabcolsep}{15pt}

\begin{tabular}{|c|c|c|}
\hline
\textbf{CO / PSO} &
\textbf{PSO1} &
\textbf{PSO2}
\\
\hline

\textbf{CO1} & & \\ \hline
\textbf{CO2} & & \\ \hline
\textbf{CO3} & & \\ \hline
\textbf{CO4} & & \\ \hline
\textbf{CO5} & & \\ \hline
\textbf{CO6} & & \\ \hline

\end{tabular}
\end{table}

% =========================================================
% CO-PO / CO-PSO MAPPING JUSTIFICATION
% =========================================================

\section*{CO--PO / CO--PSO Mapping Justification}

The CO--PO and CO--PSO mappings shall be determined based on the
actual contribution of the project towards the attainment of the
respective Course Outcomes (COs), Program Outcomes (POs), and
Program Specific Outcomes (PSOs).

For each CO, students shall identify all relevant POs and PSOs to
which the CO is mapped, specify the strength of each mapping, and
provide suitable justification based on evidence from the project.

\subsection*{Mapping Strength}

\begin{itemize}
    \item \textbf{3 -- High Contribution:} Substantial and direct
    contribution towards the attainment of the corresponding PO or PSO.

    \item \textbf{2 -- Moderate Contribution:} Meaningful but partial
    contribution towards the attainment of the corresponding PO or PSO.

    \item \textbf{1 -- Low Contribution:} Limited or indirect
    contribution towards the attainment of the corresponding PO or PSO.

    \item \textbf{-- -- No Significant Contribution:} No meaningful
    contribution towards the corresponding PO or PSO.
\end{itemize}

\vspace{0.5cm}

Students should map only those POs and PSOs for which there is
sufficient evidence from the project. The strength of each mapping
shall be assigned independently. Additional POs or PSOs may be
included wherever applicable.

% =========================================================
% CO1
% =========================================================

\subsection*{CO1}

\textbf{CO1:} Articulate the nature of science, differentiate between
objective and subjective thinking, and situate scientific inquiry within
its historical and philosophical context.

\vspace{0.3cm}

\textbf{POs mapped:} \rule{10cm}{0.4pt}

\textbf{Strength of mapping:} \rule{9cm}{0.4pt}

\textbf{Justification:}

\vspace{2cm}

\textbf{PSOs mapped:} \rule{10cm}{0.4pt}

\textbf{Strength of mapping:} \rule{9cm}{0.4pt}

\textbf{Justification:}

\vspace{2cm}


% =========================================================
% CO2
% =========================================================

\subsection*{CO2}

\textbf{CO2:} Apply inductive, deductive, and syllogistic reasoning to
construct and evaluate valid scientific arguments and identify causal
relationships.

\vspace{0.3cm}

\textbf{POs mapped:} \rule{10cm}{0.4pt}

\textbf{Strength of mapping:} \rule{9cm}{0.4pt}

\textbf{Justification:}

\vspace{2cm}

\textbf{PSOs mapped:} \rule{10cm}{0.4pt}

\textbf{Strength of mapping:} \rule{9cm}{0.4pt}

\textbf{Justification:}

\vspace{2cm}


% =========================================================
% CO3
% =========================================================

\subsection*{CO3}

\textbf{CO3:} Apply statistical concepts including the Central Limit
Theorem, error propagation, confidence intervals, and data
visualization techniques to interpret experimental results.

\vspace{0.3cm}

\textbf{POs mapped:} \rule{10cm}{0.4pt}

\textbf{Strength of mapping:} \rule{9cm}{0.4pt}

\textbf{Justification:}

\vspace{2cm}

\textbf{PSOs mapped:} \rule{10cm}{0.4pt}

\textbf{Strength of mapping:} \rule{9cm}{0.4pt}

\textbf{Justification:}

\vspace{2cm}


% =========================================================
% CO4
% =========================================================

\subsection*{CO4}

\textbf{CO4:} Design and execute hypothesis testing procedures using
Z-test, T-test, and Chi-Square tests, and critically evaluate the
validity of scientific claims.

\vspace{0.3cm}

\textbf{POs mapped:} \rule{10cm}{0.4pt}

\textbf{Strength of mapping:} \rule{9cm}{0.4pt}

\textbf{Justification:}

\vspace{2cm}

\textbf{PSOs mapped:} \rule{10cm}{0.4pt}

\textbf{Strength of mapping:} \rule{9cm}{0.4pt}

\textbf{Justification:}

\vspace{2cm}


% =========================================================
% CO5
% =========================================================

\subsection*{CO5}

\textbf{CO5:} Construct mathematical models, perform dimensional
analysis, and derive functional relationships to support theoretical
research.

\vspace{0.3cm}

\textbf{POs mapped:} \rule{10cm}{0.4pt}

\textbf{Strength of mapping:} \rule{9cm}{0.4pt}

\textbf{Justification:}

\vspace{2cm}

\textbf{PSOs mapped:} \rule{10cm}{0.4pt}

\textbf{Strength of mapping:} \rule{9cm}{0.4pt}

\textbf{Justification:}

\vspace{2cm}


% =========================================================
% CO6
% =========================================================

\subsection*{CO6}

\textbf{CO6:} Produce well-structured scientific documents, deliver
conference presentations, and demonstrate ethical conduct in research
and publication.

\vspace{0.3cm}

\textbf{POs mapped:} \rule{10cm}{0.4pt}

\textbf{Strength of mapping:} \rule{9cm}{0.4pt}

\textbf{Justification:}

\vspace{2cm}

\textbf{PSOs mapped:} \rule{10cm}{0.4pt}

\textbf{Strength of mapping:} \rule{9cm}{0.4pt}

\textbf{Justification:}

\vspace{2cm}
% =========================================================
% GUIDE VERIFICATION
% =========================================================

\clearpage
\section*{Guide Verification}

I hereby verify that the CO--PO and CO--PSO mappings of this project
have been reviewed and corrected, wherever necessary, to accurately
reflect the contribution of the project towards the prescribed Course
Outcomes, Program Outcomes, and Program Specific Outcomes.

\vspace{2.5cm}

\begin{flushright}

\begin{tabular}{c}

\rule{5cm}{0.4pt} \\[0.3cm]

\textbf{\guidename} \\[0.2cm]

\guideDesignation \\[0.3cm]

Date: \rule{3cm}{0.4pt}

\end{tabular}

\end{flushright}

% =========================================================
% J. ADDITIONAL EVIDENCE
% =========================================================

\section{Additional Evidence}

Include any other relevant evidence, certificates, photographs,
awards, presentations, technical reports, or supporting documents
associated with the major project.
